
\section{Appendix}

In the supplement, we present more results on geometry, ablation study on rendering, time analysis of our pipeline and data explanation in the following.

\begin{figure}[h!]
  \centering
  \includegraphics[width=1.\linewidth]{figs/x_comp_geo.pdf}
  \vspace{-6mm}
  \caption{\textbf{Qualitative comparison on geometry.} From left to right, we show (a) source view images, the reconstruction point maps with corresponding pixels of (b) DUSt3R, (c) Pow3R, (d) VGGT, (e) MASt3R+Prompt-DA, and (f) Ours.}
  \label{fig:x_geo}
  %\vspace{-3mm}
\end{figure}

\section{More Results on Geometry}

Our geometry module in stage 1 is able to generalize to different camera types, although our model is trained in a self-supervised manner without 3D loss.
%
In addition to the qualitative results on Mobile data in Fig.~\ref{fig:geo_ablation}, we illustrate the robustness of our geometry module in Fig.~\ref{fig:x_geo} on Industry Camera data.
%
DUSt3R~\cite{wang2024dust3r} and Pow3R~\cite{jang2025pow3r} struggle with large misalignment between two point maps on the foreground in Fig.~\ref{fig:x_geo}(b,c), even if Pow3R requires camera calibration.
%
Under only 2 sparse input views, VGGT~\cite{wang2025vggt} also badly aligns the foreground part of the point map, see face, ball and trousers in Fig.~\ref{fig:x_geo}(d).
%
By feeding MASt3R~\cite{leroy2024grounding} geometry as input, Prompt-DA~\cite{lin2025prompting} manages to refine the geometry of MASt3R but in a view-independent manner.
%
Thus the global alignment of two source views is not held.
%
Our affine transform module preserves the consistency of the overlapped part of two source views by using rendering loss and geometry regularization. 

\begin{figure}[h!]
  \centering
  \includegraphics[width=\linewidth]{figs/train_stage1_aaai65.pdf}
  \caption{\textbf{Self-supervised training of stage 1.} We train the affine transform module and auxiliary layers by using Chamfer distance and rendering loss.}
  \label{fig:stage1}
\end{figure}

\begin{table}[t!]
%\setlength\tabcolsep{2.5pt} 
\small
\centering
%\vspace{-2mm}
\begin{tabular}{l|ccc}
\toprule

\multicolumn{1}{c}{\multirow{1}[0]{*}{Method}} & PSNR$\uparrow$ & SSIM$\uparrow$ & LPIPS$\downarrow$ \\ \midrule

Stage 1 Render          & 24.844 & 0.794 & 0.296 \\
Stage 2 Initial Color          & 27.308 & 0.856 & 0.169 \\
Stage 2 Final Splatting         & \textbf{28.703} & \textbf{0.889} & \textbf{0.169} \\

\bottomrule
\end{tabular}
\caption{\textbf{Ablation study of rendering module.} We average the metrics across all datasets.}%We report Chamfer distance in centimeter for the direction of prediction to ground truth amd ground truth to prediction.}
%\vspace{-2mm}
\label{tab:render}
\end{table}

\begin{figure}[h!]
  \centering
  \includegraphics[width=1.\linewidth]{figs/render_abla4.pdf}
  \vspace{-6mm}
  \caption{\textbf{Qualitative ablation results.} We show novel view synthesis results of (a) stage 1, (b) initial warping color of stage 2, (c) final splatting of stage 2, and (d) ground truth.}
  \label{fig:render_ablation}
  %\vspace{-3mm}
\end{figure}



\section{Ablation Study on Rendering}
As mentioned in the Method section, our stage 1 network is associated with some auxiliary layers to generate two Gaussian planes, in order to supervise the geometry with rendering loss in a self-supervised manner, see Fig.~\ref{fig:stage1}.
%
However, such rendering typically struggles with some missing parts, Fig.~\ref{fig:render_ablation}(a) and obtains a low numeric result in Tab.~\ref{tab:render}, because the coarse registration in stage 1 can not totally handle the large sparsity of input views. 
%
After geometry refinement in stage 2, we can render the target view by warping the color from source views as in Eq.~\ref{eq:color_stage2}.  
%
However, the floating points between foreground and background would make noisy results, Fig.~\ref{fig:render_ablation}(b). 
%
Therefore, we use the color residual map in Eq.~\ref{eq:color} and Splatting mechanism to correct the artifacts in Fig~\ref{fig:render_ablation}(c), and to improve the quantitative result in Tab.~\ref{tab:render}. 

\begin{table}[t!]

\small
\centering
%\vspace{-2mm}
\begin{tabular}{c|l|c c}
\toprule

\multicolumn{2}{c}{\multirow{1}[0]{*}{Module}} & Time (ms) & Inp. Res.\\
\midrule
\multirow{2}[5]{*}{\rotatebox{90}{\scriptsize{Stage 1}}}
& Point Init. (MASt3R)          & 97 & \multirow{2}[5]{*}{\rotatebox{90}{\tiny{$512 \times 288$}}}  \\
& Affine Transform             & 34 & \\
& Depth Init.               & 5  & \\
\midrule
\multirow{3}[5]{*}{\rotatebox{90}{\scriptsize{Stage 2}}}& Depth Refine.             & 119 & \multirow{3}[5]{*}{\rotatebox{90}{\scriptsize{$1024 \times 576$}}} \\
& Color Init.                  & 5  & \\
& Gaussian Plane/Color Correct & 18 &  \\
& Splatting                    & 2  &  \\
\midrule
& Total      & 280 & \\

\bottomrule
\end{tabular}
\caption{\textbf{Time cost of our pipeline.} In stage 1, our network takes a pair of images in the resolution of $512 \times 288$ as inputs, while two images of $1024 \times 576$ are fed into our refinement module in stage 2. Our full pipeline takes totally 280ms.}
%\vspace{-2mm}
\label{tab:time_analysis}
\end{table}

\section{Time Analysis}
We conduct experiments on a machine equipped with an RTX 3090 GPU with 24GB memory for our method and provide a time analysis of our pipeline in Tab.~\ref{tab:time_analysis}.
%
In stage 1, point map reconstruction takes a lot of time by using original MASt3R~\cite{leroy2024grounding}, due to the complex structure of ViT, while our iterative affine transform is very efficient.  
%
In addition, the depth initialization is achieved by using $\alpha$-blending~\cite{kerbl2023_3dgs} with a fixed radius and diagonal rotation matrix.
%
Due to the fine resolution input in stage 2, the majority of time is used for depth refinement by using the costly 3D convolution.
%
The time analysis is done by using PyTorch and we believe that the full pipeline can be largely accelerated with a C++ implementation of TensorRT.

\section{Data} 

As mentioned in the Experiment section, we train and evaluate our method on 3 types of camera data, including industry camera (THumanMV~\cite{zhou2024gps}), GoPro, and mobile phone (4K4D~\cite{xu20244k4d} and SelfCap~\cite{xu2024representing}). 
%
For industry camera, we take 15/11 training/validation sequences in different scenes.
The validation data is, in general, unseen character or unseen motion.
For mobile stage data, we take around 3000 frames from 4K4D dance sequence and SelfCap yoga sequence for training and validate on the rest of the frames.
Since the industry camera and the mobile phone stage are typically complex capture systems with relatively large focal lengths, they can only capture small amplitude movement.
%
Therefore, we capture multi-person movement of sport in large scale scenes with a portable GoPro system in the mode of 1080P 30FPS.
%
For GoPro data, we take 6/4 sequences as training/validation data.
%
All aforementioned datasets provide camera calibration.
%
The majority of the data we used has been already publicly available for research purposes, while a part of the data is not available due to a confidential issue.


